\documentclass[journal]{IEEEtran}
\usepackage{graphicx}

\ifCLASSOPTIONcompsoc
\usepackage[caption=false,font=normalsize,labelfont=sf,textfont=sf]{subfig}
\else
\usepackage[caption=false,font=footnotesize]{subfig}
\fi

\hyphenation{op-tical net-works semi-conduc-tor}

\begin{document}

\title{Distal Femur Segmentation\\on MR Images Using Random Forests}

\author{Fabian~Balsiger, %~\IEEEmembership{Member,~IEEE,}
        Tiziano~Ronchetti, %~\IEEEmembership{Fellow,~OSA,}
        and~Martin~Pletscher %,~\IEEEmembership{Life~Fellow,~IEEE}% <-this % stops a space
\thanks{The authors are students at the department of Biomedical Engineering, University of Bern, Switzerland. Use e-mails fabian.balsiger@students.unibe.ch and martin.pletscher@students.unibe.ch to contact F. Balsiger and M. Pletscher, respectively. T. Ronchetti can be contacted under tiziano.ronchetti@bfh.ch.}% <-this % stops a space
\thanks{Manuscript received January 7, 2017.}}

\markboth{Medical Image Analysis Laboratory 2017}%
{Shell \MakeLowercase{\textit{et al.}}: MIA Lab}

\maketitle

% As a general rule, do not put math, special symbols or citations
% in the abstract or keywords.
\begin{abstract}
Segmentation of the distal femur from magnetic resonance (MR) images plays a crucial role in clinical practice. Diagnosis, disease progression monitoring, preoperative planning, and treatment analysis of diseases like osteoarthritis are applications. We developed a fully automatic method to segment the distal femur on MR images. Our method uses random forests (RF) and consists of two main phases: training and segmentation, i.e., a RF model is trained that enables RF voting to classify the distal femur on MR images. We evaluated our method with 20 MR data sets manually labeled by experts. A leave-one-out cross-validation was performed ($n=10$) and Dice similarity coefficients (DSC) were obtained. The mean DSC is $92.37 \pm 0.0266$. Mean sensitivity is $91.57 \%$ and mean specificity is $99.29 \%$. Training the RF model with 19 data sets took in average 791 seconds and segmenting an image 45 seconds ($n=200$). Our method is suitable to segment the distal femur with high accuracy and precision. Moreover, the duration to segment the femur with less than a minute is feasible for clinical use.
\end{abstract}

% Note that keywords are not normally used for peerreview papers.
\begin{IEEEkeywords}
Fully automated femur segmentation, femur detection, segmentation, random forest (RF), magnetic resonance imaging (MRI).
\end{IEEEkeywords}

\IEEEpeerreviewmaketitle

\section{Introduction}
% The very first letter is a 2 line initial drop letter followed
% by the rest of the first word in caps.
% 
% form to use if the first word consists of a single letter:
% \IEEEPARstart{A}{demo} file is ....
% 
% Here we have the typical use of a "T" for an initial drop letter
% and "HIS" in caps to complete the first word.
\IEEEPARstart{T}{his} demo file is intended to serve as a ``starter file''
for IEEE journal papers produced under \LaTeX\ using
IEEEtran.cls version 1.8b and later.
% You must have at least 2 lines in the paragraph with the drop letter
% (should never be an issue)
We wish you the best of success.


\subsection{Subsection Heading Here}
Subsection text here. This is a citation~\cite{Criminisi2013}.

% needed in second column of first page if using \IEEEpubid
%\IEEEpubidadjcol

\subsubsection{Subsubsection Heading Here}
Subsubsection text here.


% An example of a floating figure using the graphicx package.
% Note that \label must occur AFTER (or within) \caption.
% For figures, \caption should occur after the \includegraphics.
% Note that IEEEtran v1.7 and later has special internal code that
% is designed to preserve the operation of \label within \caption
% even when the captionsoff option is in effect. However, because
% of issues like this, it may be the safest practice to put all your
% \label just after \caption rather than within \caption{}.
%
% Reminder: the "draftcls" or "draftclsnofoot", not "draft", class
% option should be used if it is desired that the figures are to be
% displayed while in draft mode.
%
\begin{figure}[!h]
	\centering
	\includegraphics[width=2.5in]{logo}
	% where an .eps filename suffix will be assumed under latex, 
	% and a .pdf suffix will be assumed for pdflatex; or what has been declared
	% via \DeclareGraphicsExtensions.
	\caption{The IEEE logo.}
	\label{fig_logo}
\end{figure}

% Note that the IEEE typically puts floats only at the top, even when this
% results in a large percentage of a column being occupied by floats.


% An example of a double column floating figure using two subfigures.
% (The subfig.sty package must be loaded for this to work.)
% The subfigure \label commands are set within each subfloat command,
% and the \label for the overall figure must come after \caption.
% \hfil is used as a separator to get equal spacing.
% Watch out that the combined width of all the subfigures on a 
% line do not exceed the text width or a line break will occur.
%
%\begin{figure*}[!t]
%\centering
%\subfloat[Case I]{\includegraphics[width=2.5in]{box}%
%\label{fig_first_case}}
%\hfil
%\subfloat[Case II]{\includegraphics[width=2.5in]{box}%
%\label{fig_second_case}}
%\caption{Simulation results for the network.}
%\label{fig_sim}
%\end{figure*}
%
% Note that often IEEE papers with subfigures do not employ subfigure
% captions (using the optional argument to \subfloat[]), but instead will
% reference/describe all of them (a), (b), etc., within the main caption.
% Be aware that for subfig.sty to generate the (a), (b), etc., subfigure
% labels, the optional argument to \subfloat must be present. If a
% subcaption is not desired, just leave its contents blank,
% e.g., \subfloat[].


% An example of a floating table. Note that, for IEEE style tables, the
% \caption command should come BEFORE the table and, given that table
% captions serve much like titles, are usually capitalized except for words
% such as a, an, and, as, at, but, by, for, in, nor, of, on, or, the, to
% and up, which are usually not capitalized unless they are the first or
% last word of the caption. Table text will default to \footnotesize as
% the IEEE normally uses this smaller font for tables.
% The \label must come after \caption as always.
%
%\begin{table}[!t]
%% increase table row spacing, adjust to taste
%\renewcommand{\arraystretch}{1.3}
% if using array.sty, it might be a good idea to tweak the value of
% \extrarowheight as needed to properly center the text within the cells
%\caption{An Example of a Table}
%\label{table_example}
%\centering
%% Some packages, such as MDW tools, offer better commands for making tables
%% than the plain LaTeX2e tabular which is used here.
%\begin{tabular}{|c||c|}
%\hline
%One & Two\\
%\hline
%Three & Four\\
%\hline
%\end{tabular}
%\end{table}


% Note that the IEEE does not put floats in the very first column
% - or typically anywhere on the first page for that matter. Also,
% in-text middle ("here") positioning is typically not used, but it
% is allowed and encouraged for Computer Society conferences (but
% not Computer Society journals). Most IEEE journals/conferences use
% top floats exclusively. 
% Note that, LaTeX2e, unlike IEEE journals/conferences, places
% footnotes above bottom floats. This can be corrected via the
% \fnbelowfloat command of the stfloats package.

\section{Materials and Methods}

\section{Results}

\section{Discussion}

\section{Conclusion}


\appendices
\section{Choice of Parameters}
\label{secAppendixA}
The method is generally insensitive to most of the parameter settings. A change in the below introduced parameters does not affect the segmentation outcome significantly. However, it can affect the computational effort to build the RF model and to segment an image.
\subsection{Random Forest Parameters}

\textit{Number of trees}: The number of trees to train the RF has been determined by trial and error to be 20. Underfitting occurred with a number of trees lower than 20 in certain cases. More than 20 trees does not improve the DSC significantly and thus has no benefit.

\textit{Number of voxels}: We used $5 \%$ of the voxels that show femur and $5 \%$ of the non-femur voxels from each data set to train the RF. Tests showed that the DSC does not improve significantly with increasing number of voxels used to train the RF. Though, the computational time is increased by factor 2.5 for each additional $5 \%$ voxels.

\textit{Volumetric neighborhood}: We determined a neighborhood of $I~x~J~x~K=5~x~5~x~3$ voxels for the volumetric features (i.e., mean, variance, and entropy) by trial and error. Tests showed that a smaller or larger neighborhood does negatively affect the RF voting outcome. A change in the neighborhood size does not significantly affect the computational effort.

\subsection{Segmentation Parameters}
\textit{Scaling factor}: The image that is segmented was scaled to a smaller size to achieve a faster segmentation. A scaling by $0.5$ has been determined for our evaluation. Segmentation without scaling does not improve the segmentation result significantly but does increase the computational time by factor 4.

\section*{Acknowledgment}
The authors would like to thank Prof. Dr. Mauricio Reyes for his generous support and helpful feedback during the laboratory and for providing the materials. He gave us an excellent introduction to the field of random forests, statistical shape models, and active shape models.

% Can use something like this to put references on a page
% by themselves when using endfloat and the captionsoff option.
\ifCLASSOPTIONcaptionsoff
  \newpage
\fi

\input{bib}

\end{document}

